\section{What an SVD Actually Contains}
\label{sec:svd-what}

An SVD (System View Description) is an XML file — CMSIS-SVD, a standard — that describes a
microcontroller's memory-mapped surface. It exists because every debugger needs the same information
and nobody should re-derive it per vendor.

\subsection{The Five Things That Matter}

\begin{table}[htbp]
    \centering
    \footnotesize
    \begin{tabular}{@{}p{0.20\textwidth}p{0.34\textwidth}p{0.42\textwidth}@{}}
        \toprule
        \textbf{Element} & \textbf{What it says} & \textbf{Where you meet it in Rust} \\
        \midrule
        \codeinline{peripheral} & base address and register layout & \codeinline{Periph<T, ADDR>} \\
        \codeinline{register} & offset, size, access (r/w/rw) & \codeinline{registers.read()}, \codeinline{.write()} \\
        \codeinline{field} & bit offset, width, access & \codeinline{w.mvp().set\_bit()} \\
        \codeinline{enumeratedValue} & named values, and their descriptions & the enum in the field writer \\
        \codeinline{resetValue} & what the hardware holds after reset & \codeinline{reset()} \\
        \bottomrule
    \end{tabular}
    \caption{The five elements that produce all the code you will actually call. Everything else in
    an SVD — interrupts, dim arrays, clusters — is scaffolding around them.}
    \label{tab:svd-elements}
\end{table}

\subsection{A Register, in the SVD and in Rust}

The MVP's command register, as the SVD describes it:

\begin{minted}{text}
<register>
  <name>CMD</name>
  <addressOffset>0x174</addressOffset>
  <fields>
    <field><name>START</name><bitOffset>0</bitOffset><bitWidth>1</bitWidth>
           <description>Start Command</description></field>
    <field><name>INIT</name><bitOffset>3</bitOffset><bitWidth>1</bitWidth>
           <description>Initialization Command/Qualifier</description></field>
    ...
\end{minted}

and what that becomes, which is what you type:

\begin{minted}{rust}
// One write, two fields: the qualifier and the command together.
mvp.cmd().write(|w| w.init().set_bit().start().set_bit());
\end{minted}

The word \emph{Qualifier} in the description is the whole interface contract for that register: it
says INIT modifies the meaning of the command in the same write rather than being a separate command
of its own. Reading the generated code is reading the SVD.

\subsection{Where the SVD Comes From, and Why It Is Trustworthy}

Silicon Labs publishes SVD files per part, and their own tools build register views from the same
file. That is the reason to trust it more than a PDF: a manual is written once and then slowly drifts,
while an SVD that disagrees with the hardware breaks the vendor's debugger on the same day.

It is still not a specification of the \emph{board}. Three of the errors documented in
section~\ref{sec:gpio-vendor-errors} are board facts — which pin the LED is on, which way the I2C
lines go — and no SVD contains them. Chapter~4's rule applies: chip facts from the SVD, board facts
from a measurement.

\subsection{The Practical Consequence: One File, One Crate}

The project's PAC is generated once and checked in, because a build that downloads an SVD from the
network is a build that stops working when the network does. Regeneration is a deliberate act:

\begin{minted}{bash}
# the SVD lives beside the crate it generates
ls vendor/EFR32MG24B220F1536IM48.svd
svd2rust -i vendor/EFR32MG24B220F1536IM48.svd --target cortex-m
# then split the generated lib.rs into one module per peripheral, which is what
# makes `cargo doc` usable and keeps a 40k-line file reviewable
\end{minted}

That split is why this repository's PAC is a directory of small modules rather than one enormous
file: \codeinline{mvp\_s/}, \codeinline{gpio\_s/}, \codeinline{cmu\_s/}, and so on. It costs a script
and saves every future diff.
